\section{Related Work}
\label{sec:related}

The economics of passive provision is now well understood. Milionis et al.\
formalised the loss a constant-product LP suffers against a delta-rebalanced
portfolio and derived the rate $\frac{1}{8}\sigma^2 V$ for a full-range
position; the loss appears on any price path, informed traders or not, and
protocol fees and the choice of curve shift only the balance\cite{milionis2022lvr}.
Fritsch and Canidio measured the losses on real pools and found that arbitrage
losses typically dominate the fees earned, with passive full-range positions
outperforming concentrated ones in their sample\cite{fritsch2024measuring}.
Loesch et al.\ studied v3 positions directly and showed that concentration
magnifies impermanent loss when the price leaves the band\cite{loesch2021impermanent}.
Together these results define the terrain: gross fee yield is not an LP's
return.

A second line of work learns the LP's decisions. Zhang et al.\ train a hedged
double deep Q-network and report large gains over baselines, but the hedge is
assumed frictionless and the policy sees no regime
features\cite{zhang2023adaptive}. A recent replication replaces the feature
extractor of that model with a sequence model and a revised reward, and adds
two baselines; it keeps the simulated setting and reports no statistic that
adjusts for risk\cite{zhang2025mamba}. ZeroSwap learns the spread of a market maker so
that quoted prices track an external reference, an orthogonal lever to range
placement\cite{nadkarni2023zeroswap}. RAmmStein poses regime-dependent
rebalancing thresholds as a problem of impulse control and reports that
rebalancing lazily when the regime warrants it improves net ROI with far fewer
rebalances, though its net metric excludes impermanent loss and the study
relies on off-chain data for one pair\cite{anchuri2026rammstein}. Chionas et
al.\ treat concentrated liquidity as a risk-aversion problem in simulation and
find that higher volatility and gas widen optimal positions and reduce
rebalancing, while a risk penalty trims the adverse tail; their analysis uses
no on-chain data and no hedge\cite{chionas2026concentrated}. Classical market
making supplies the risk-sensitive vocabulary we use, but its state is a limit
order book and its actions are quotes, not range geometry\cite{avellaneda2008high}.

The closest study to this report trains a PPO agent on hourly WETH/USDC
$0.05\,\%$ data from the same primary pool used here\cite{xu2025defi}. Xu and
Brini let the policy decide when to redeploy and how wide to keep the band,
penalise the loss-versus-rebalancing rate inside the reward, and compare the
learned policy with a periodic reset across eleven rolling windows; the active
policy wins seven of them. Three differences matter for the claims of this
report. Their agent carries no hedge, so directional risk stays in the
position. Market outcomes are simulated from closed-form fee and loss
expressions at hourly resolution, with a fixed gas charge and no slippage,
rather than replayed from the recorded events. And the evaluation is thinner:
performance is reported as cumulative reward, without a statistic that adjusts
for risk or an uncertainty interval, and the best of fifty trained agents is
selected for each window.

The learning machinery is standard. We use proximal policy optimisation with
clipping (PPO)\cite{schulman2017ppo}, implemented through Stable-Baselines3 on
Gymnasium environments\cite{raffin2021stablebaselines3,towers2024gymnasium},
and evaluate risk with CVaR at the $5\,\%$ level through its
Rockafellar--Uryasev representation\cite{rockafellar2000cvar}.

No study surveyed here combines the parts this report tests. The line that
learns decisions supplies them only in pieces: a hedge learned without
friction, a range policy learned without the hedge, and a mean metric rather
than one that adjusts for risk. Two of the answers here are measured rather
than assumed. The passive benchmarks are decomposed into fees, impermanent
loss, gas and slippage on both fee tiers, so the economics answer rests on the
pools of this study rather than on a survey, and the ways a standard policy
collapses on the task are reported as a result rather than hidden in an
implementation note.

What remains is the combination studied here. One risk-sensitive policy
chooses timing, width and hedge. A cost ledger replayed from the tape,
bounded against an exact integer engine, prices every friction on both legs.
Sharpe and tail statistics carry paired intervals across seeds. And a boundary
states where active management stops paying. The rest of the report tests that
combination, with the hedge leg the one part still pending: the overlay is
disabled in the shipped configuration, so the hedge results are marked as such.
